%% kapa: exam=1 round=21 year=2010 period=2 form=A subject=accounting subject_ko=회계학 seq=4 number=4
%% source: Q-Net 감정평가사 기출문제 (cst00309) artlSeq=5024181 file=extracted/1008/5024181/2010년 제21회 감정평가사 1차시험 기출문제/2교시.hwp
%% conversion: kordoc 4.24.51 -> tools/gen1.py (auto); review=proofread-sonnet; flags=none
\begin{question}{4}
20×1년 초 (주)한국은 (주)서울의 지분 100\%를 취득하여 흡수합병하였고, 그 이전대가로 ￦100,000의 현금을 지급하였다. 또한, 합병관련 컨설팅 수수료로 ￦10,000을 현금으로 지급하였다. 합병 당시 (주)서울의 재무상태표는 다음과 같다.

\begin{center}\uline{재무상태표}\end{center}
\noindent (주)서울 \hfill 20×1년 1월 1일 현재 \hfill (단위: 원)\par

{\small\centering
\begin{tabular}{l r | l r}
\hline
매출채권 & ￦24,000 & 매입채무 & ￦48,000 \\
상품 & 30,000 & 납입자본 & 40,000 \\
토지 & \uline{54,000} & 이익잉여금 & \uline{20,000} \\
자산총계 & ￦108,000 & 부채와 자본총계 & ￦108,000 \\
\end{tabular}\par}

20×1년 초 (주)한국이 (주)서울의 자산에 대해 공정가치를 측정한 결과, 매출채권은 장부금액과 동일하며 상품은 장부금액 대비 10\% 높고 토지는 장부금액 대비 50\% 높은 것으로 평가되었다. 한편, 부채의 공정가치는 장부금액과 동일하다. 20×1년 초 (주)한국이 영업권으로 계상할 금액은 얼마인가?
\begin{choices}
  \item ￦5,000
  \item ￦10,000
  \item ￦15,000
  \item ￦20,000
  \item ￦23,000
\end{choices}
\end{question}
\answer{2}
